% ============================================================
%  IntelIPath — Project Report
%  Compile with: pdflatex / xelatex
% ============================================================
\documentclass[12pt,a4paper]{report}

% ---------- Packages ----------
\usepackage[utf8]{inputenc}
\usepackage[T1]{fontenc}
\usepackage[english]{babel}
\usepackage[a4paper,margin=2.5cm]{geometry}
\usepackage{graphicx}
\usepackage{float}
\usepackage{caption}
\usepackage{subcaption}
\usepackage{booktabs}
\usepackage{longtable}
\usepackage{array}
\usepackage{enumitem}
\usepackage{xcolor}
\usepackage{hyperref}
\usepackage{fancyhdr}
\usepackage{titlesec}
\usepackage{tocloft}
\usepackage{amsmath}
\usepackage{tabularx}

% ---------- Hyperlink style ----------
\hypersetup{
    colorlinks=true,
    linkcolor=black,
    citecolor=black,
    urlcolor=blue,
    pdftitle={IntelIPath - Project Report},
    pdfauthor={IntelIPath Team}
}

% ---------- Header / Footer ----------
\pagestyle{fancy}
\fancyhf{}
\fancyhead[L]{IntelIPath}
\fancyhead[R]{\leftmark}
\fancyfoot[C]{\thepage}
\renewcommand{\headrulewidth}{0.4pt}

% ---------- Section numbering depth ----------
\setcounter{secnumdepth}{3}
\setcounter{tocdepth}{3}

% ---------- Helper: image placeholder box ----------
% Usage: \figplaceholder{width}{height}{Caption text}{label}
\newcommand{\figplaceholder}[4]{%
    \begin{figure}[H]
        \centering
        \fbox{%
            \begin{minipage}[c][#2][c]{#1}
                \centering
                \textit{[Insert image here: #3]}
            \end{minipage}%
        }
        \caption{#3}
        \label{#4}
    \end{figure}
}

% Real image insertion (uncomment and edit path when the image is ready):
% \begin{figure}[H]
%     \centering
%     \includegraphics[width=0.9\textwidth]{images/image-name.png}
%     \caption{Figure description}
%     \label{fig:image-name}
% \end{figure}

\begin{document}

% ============================================================
% COVER PAGE
% ============================================================
\begin{titlepage}
    \centering
    \includegraphics[width=0.3\textwidth]{images/school-logo.png}\\[0.5cm]
    % TODO: replace with the school/project logo at images/school-logo.png

    {\large \textbf{FPT UNIVERSITY}}\\[0.2cm]
    {\normalsize DEPARTMENT OF INFORMATION TECHNOLOGY}\\[2.5cm]

    {\Huge \textbf{PROJECT REPORT}}\\[0.4cm]
    {\Large \textbf{IntelIPath — Personalized Learning Roadmap Platform}}\\[2cm]

    \begin{tabular}{ll}
        \textbf{Course:}      & \textit{[Course name]} \\
        \textbf{Supervisor:}  & \textit{[Supervisor name]} \\
        \textbf{Class:}       & \textit{[Class code]} \\
        \textbf{Semester:}    & \textit{[Semester / Year]} \\
    \end{tabular}

    \vfill

    {\large \today}
\end{titlepage}

% ============================================================
% TABLE OF CONTENTS
% ============================================================
\tableofcontents
\listoffigures
\listoftables
\clearpage

% ============================================================
% CHAPTER: INTRODUCTION
% ============================================================
\chapter{Introduction}

\section{Overview}
IntelIPath is a platform that helps students build and track a personalized learning roadmap,
connecting students with mentors and academic counselors to receive feedback on their
portfolio and career guidance. The system is role-based (Student, Mentor, Counselor, Admin)
and integrates a Virtual Mentor assistant for automated guidance.

% TODO: expand this section — project scope, target users, core technology stack
% (Spring Boot backend, React + TypeScript frontend, PostgreSQL, Docker).

\section{Project Objectives}
\begin{itemize}[leftmargin=1.2cm]
    \item Build a visual learning roadmap organized by skill/career track.
    \item Allow students to build a portfolio and request reviews from mentors.
    \item Support academic counselors in tracking and managing student progress.
    \item \textit{[Add further objectives if applicable]}
\end{itemize}

% ============================================================
% CHAPTER: MEMBER / TEAM
% ============================================================
\chapter{Team Members}

\begin{table}[H]
    \centering
    \caption{Team members and roles}
    \label{tab:members}
    \begin{tabularx}{\textwidth}{|l|X|X|X|}
        \hline
        \textbf{Student ID} & \textbf{Full Name} & \textbf{Role} & \textbf{Main Responsibility} \\
        \hline
        \textit{[Student ID]} & \textit{[Full name]} & \textit{[Team Lead / Backend / Frontend / ...]} & \textit{[Responsibility]} \\
        \hline
        \textit{[Student ID]} & \textit{[Full name]} & \textit{[Role]} & \textit{[Responsibility]} \\
        \hline
        \textit{[Student ID]} & \textit{[Full name]} & \textit{[Role]} & \textit{[Responsibility]} \\
        \hline
        \textit{[Student ID]} & \textit{[Full name]} & \textit{[Role]} & \textit{[Responsibility]} \\
        \hline
    \end{tabularx}
\end{table}

% ============================================================
% CHAPTER: PROBLEMS
% ============================================================
\chapter{Problems}

\section{Background}
% TODO: describe the current state before the system existed — students struggle to build
% their own learning roadmap, lack a direct channel to mentors, cannot centrally track
% progress, etc.

\section{Problem Statement}
\begin{enumerate}[leftmargin=1.2cm]
    \item \textit{[Problem 1: e.g. students lack a clear roadmap organized by skill]}
    \item \textit{[Problem 2: e.g. no feedback channel from mentors/counselors]}
    \item \textit{[Problem 3: e.g. counselors struggle to manage a large number of students]}
    \item \textit{[Problem 4]}
\end{enumerate}

\section{Affected Stakeholders}
\begin{itemize}[leftmargin=1.2cm]
    \item \textbf{Students:} \textit{[describe impact]}
    \item \textbf{Mentors:} \textit{[describe impact]}
    \item \textbf{Academic Counselors:} \textit{[describe impact]}
\end{itemize}

% ============================================================
% CHAPTER: SOLUTION
% ============================================================
\chapter{Solution}

\section{Proposed Solution}
% TODO: present the overall solution — a web platform that lets students build a learning
% roadmap, build a portfolio, and request reviews; mentors approve and give feedback;
% counselors track progress.

\section{Key Features}
\begin{table}[H]
    \centering
    \caption{Feature list by role}
    \label{tab:features}
    \begin{tabularx}{\textwidth}{|l|X|}
        \hline
        \textbf{Role} & \textbf{Feature} \\
        \hline
        Student & View learning roadmap, build portfolio, request portfolio review, reset password via email \\
        \hline
        Mentor & View/approve review requests, give feedback on student portfolios \\
        \hline
        Counselor & Track student progress, export reports, manage assigned student list \\
        \hline
        Admin & Manage users, configure the system \\
        \hline
        \textit{[...]} & \textit{[...]} \\
        \hline
    \end{tabularx}
\end{table}

\section{Technology Architecture}
% TODO: list the technology stack — Backend: Spring Boot, PostgreSQL, Docker; Frontend:
% React, TypeScript, Vite, Tailwind; Auth: JWT + OAuth2 (GitHub/Google) + magic-link reset.

\figplaceholder{0.85\textwidth}{6cm}{System Architecture Diagram}{fig:architecture}

% ============================================================
% CHAPTER: BUSINESS PROCESS
% ============================================================
\chapter{Business Process}

\section{Sign-up / Sign-in Process}
% TODO: describe the steps — sign in via OAuth (GitHub/Google) or username/password (FPT
% students), forgot password via a magic link sent by email.

\figplaceholder{\textwidth}{7cm}{Flowchart — Sign-in / Password Reset Process}{fig:process-login}

\section{Portfolio Review Request Process}
% TODO: describe the steps — student enters mentor's email, system creates the request,
% mentor gets notified, mentor approves/rejects, mentor sends feedback, student receives it.

\figplaceholder{\textwidth}{7cm}{Flowchart — Portfolio Review Request Process}{fig:process-review}

\section{Counselor Progress Tracking Process}
% TODO: describe the steps — counselor views the student list, views statistics, exports
% an Excel report.

\figplaceholder{\textwidth}{7cm}{Flowchart — Student Progress Tracking Process}{fig:process-counselor}

% ============================================================
% CHAPTER: BUSINESS RULES
% ============================================================
\chapter{Business Rules}

\begin{table}[H]
    \centering
    \caption{Business rules list}
    \label{tab:business-rules}
    \begin{tabularx}{\textwidth}{|c|X|X|}
        \hline
        \textbf{ID} & \textbf{Rule Description} & \textbf{Applies To} \\
        \hline
        BR-01 & A student may have at most \textbf{one} pending review request to the same mentor at a time & Student, Portfolio Review \\
        \hline
        BR-02 & A password-reset token (magic link) can be used \textbf{only once} and expires after a fixed period & Auth, Password Reset \\
        \hline
        BR-03 & An account must be in \textit{ACTIVE} status to sign in & Auth \\
        \hline
        BR-04 & \textit{[Add another rule]} & \textit{[...]} \\
        \hline
        BR-05 & \textit{[Add another rule]} & \textit{[...]} \\
        \hline
    \end{tabularx}
\end{table}

% ============================================================
% CHAPTER: STATE CHART
% ============================================================
\chapter{State Chart}

\section{Portfolio Review Request State Chart}
% TODO: describe the states: PENDING -> APPROVED / REJECTED -> COMPLETED, with the
% transition conditions.

\figplaceholder{0.8\textwidth}{6cm}{State chart — Portfolio review request lifecycle}{fig:state-review}

\section{User Account State Chart}
% TODO: describe the states: PENDING_ACTIVATION -> ACTIVE -> DISABLED, etc.

\figplaceholder{0.8\textwidth}{6cm}{State chart — User account lifecycle}{fig:state-account}

% ============================================================
% CHAPTER: CONTEXT DIAGRAM
% ============================================================
\chapter{Context Diagram}

% TODO: describe the system as a black box and the external actors interacting with it:
% Student, Mentor, Counselor, Admin, Email Service, OAuth Provider (GitHub/Google).

\figplaceholder{0.9\textwidth}{8cm}{Context Diagram — IntelIPath System}{fig:context-diagram}

% ============================================================
% CHAPTER: CONCEPTUAL MODEL
% ============================================================
\chapter{Conceptual Model}

\section{Key Entities}
\begin{itemize}[leftmargin=1.2cm]
    \item \textbf{User} — the shared user account (email, username, password\_hash, role, status)
    \item \textbf{Student} — the student profile, linked to \textit{User}
    \item \textbf{AcademicCounselor} — the academic counselor profile
    \item \textbf{Mentor} — the mentor profile
    \item \textbf{Portfolio} — the student's competency portfolio
    \item \textbf{PortfolioReviewRequest} — a review request sent to a mentor
    \item \textbf{Roadmap} — the learning roadmap organized by skill/career track
    \item \textit{[Add further entities]}
\end{itemize}

\figplaceholder{0.9\textwidth}{8cm}{Conceptual Model — Relationships among key entities}{fig:conceptual-model}

% ============================================================
% CHAPTER: ERD
% ============================================================
\chapter{Entity Relationship Diagram (ERD)}

% TODO: insert a detailed ERD (tables, primary/foreign keys, data types) generated from
% docker/postgres/init/01_init_intelipath.sql or a tool such as dbdiagram/drawio.

\figplaceholder{\textwidth}{10cm}{Entity Relationship Diagram (ERD) — IntelIPath Database}{fig:erd}

\section{Key Table Descriptions}
\begin{table}[H]
    \centering
    \caption{Description of key data tables}
    \label{tab:erd-tables}
    \begin{tabularx}{\textwidth}{|l|X|}
        \hline
        \textbf{Table} & \textbf{Description} \\
        \hline
        users & User account information \\
        \hline
        students & Student profile \\
        \hline
        academic\_counselor & Academic counselor profile \\
        \hline
        password\_reset\_tokens & Password-reset tokens (magic link) \\
        \hline
        portfolio\_review\_requests & Portfolio review requests \\
        \hline
        \textit{[...]} & \textit{[...]} \\
        \hline
    \end{tabularx}
\end{table}

% ============================================================
% CHAPTER: SWIMLANE
% ============================================================
\chapter{Swimlane Diagram}

\section{Swimlane — Review Request and Feedback Process}
% TODO: describe the lanes: Student | System | Mentor, showing the flow across each role.

\figplaceholder{\textwidth}{9cm}{Swimlane Diagram — Portfolio Review Request Process}{fig:swimlane-review}

\section{Swimlane — Password Reset Process}
% TODO: describe the lanes: Student | System | Email Service.

\figplaceholder{\textwidth}{9cm}{Swimlane Diagram — Password Reset Process (magic link)}{fig:swimlane-reset}

% ============================================================
% APPENDIX (optional)
% ============================================================
\appendix
\chapter{Appendix}
\section{References}
\begin{itemize}[leftmargin=1.2cm]
    \item \textit{[Reference 1]}
    \item \textit{[Reference 2]}
\end{itemize}

\end{document}
